\sectionkicker{June chronology}
\subsection*{6月の変化を時系列で束ねる}
\noindent 各章を横断して主要な転換点だけを日付順に置き直す。これは因果関係を示す年表ではなく、モデル、Agent runtime、serving、安全境界、科学workflowの変化が同じ月にどのように重なったかを読むための索引である。
\medskip
\begin{center}
\scriptsize
\renewcommand{\arraystretch}{1.08}
\begin{tabularx}{\linewidth}{@{}p{0.10\linewidth}p{0.20\linewidth}X@{}}
\toprule
日付 & 層 & 観測点 \\
\midrule
06-01 & Frontier Models & MiniMaxがM3を公開し、open weights、画像・動画入力、desktop operationを発表した。 \cite{src-d8e13386974c} \\
06-04 & Agents \& Memory & OpenAIがChatGPT Dreamingのmemory architectureを発表し、persistent memoryをproduct-levelの設計対象として前面化した。 \cite{src-c3ab33c62226} \\
06-17 & AI for Science & OpenAIとMolecule.oneがnear-autonomous chemistry workflowを報告した。human selection、lab operation、validationは残されている。 \cite{src-5ff2f715f1e6} \\
06-19 & Agents \& Memory & multi-call / tool workload向けagentic scheduling研究が公開され、長時間Agentの実行順序をruntime側の設計問題として扱った。 \cite{src-e8d9fd84bddc} \\
06-24 & Inference \& Serving & custom inference chipのengineering sampleとFlashInferのBlackwell/GQA decode等のupdateが同日に現れ、hardwareからkernelまでの更新が並行した。 \cite{src-855db08b75f6,src-272a6dac5bcc} \\
06-26 & Models / Serving & GPT-5.6 seriesのlimited previewとSGLang v0.5.14が同日。frontier modelの提供状態とserving側のsupportを分けて追う必要がある。 \cite{src-30f48121decc,src-538cecdad845} \\
06-27 & Agent Safety & capability exposureとper-call authorizationを同一視できないことを、agent frameworkの実装監査から提起した研究が公開された。 \cite{src-02da0db00c2d} \\
06-29 & Serving / Safety / Paper & vLLM v0.24.0、persistent-agent security評価、long-context search failureの研究が月末に集中し、実行系・安全性・評価の更新が重なった。 \cite{src-6cc910fb0183,src-6da0c5a88fa4,src-8603fd60d637} \\
06-30 & Frontier Models & AnthropicのJune technical releasesとGemini APIのlifecycle updateが月末を締め、model/APIの提供面が引き続き動いた。 \cite{src-f6aa679babd7,src-a24fd97eb3d0} \\
\bottomrule
\end{tabularx}
\renewcommand{\arraystretch}{1.0}
\end{center}
\scriptsize
\noindent\textit{日付は各Evidence cardで確認したfirst-announced chronologyに従う。複数資料を同じ行へ置くことは同日性または近接した観測を示すだけで、資料間の因果関係や性能優位を意味しない。}
\normalsize
